% Source-aligned English for printed p.135.
\SourcePage{135}

\LayerTarget{tohoku:prop:1.9.2}
\textsc{Proposition 1.9.2.} \textit{Let $\Sigma$ be a diagram
scheme with commutation relations \textup{(}\textit{cf.} \textup{1.6)}, let $\mathbf{C}$ be an abelian category,
and let $(U_i)_{i\in I}$ be a family of generators of $\mathbf{C}$. Assume that the
origin and endpoint of every arrow of $\Sigma$ are distinct, and that no
identity morphisms $e_s$, where $s$ is a vertex, occur in the commutation
relations}\textsuperscript{2)}. \textit{Then for every $A\in \mathbf{C}$ and every
vertex $s$ of the scheme, the diagram $\varepsilon_s(A)$ whose value is $A$ at
$s$ and $0$ at every other vertex, and whose value on every arrow is $0$,
belongs to $\mathbf{C}^\Sigma$. Moreover, the family of the $\varepsilon_s(U_i)$, with
$s$ and $i$ variable, is a system of generators of $\mathbf{C}^\Sigma$.}

\LayerTarget{tohoku:prop:1.9.2:proof}
The verification is immediate. For the last assertion it suffices to observe
that morphisms from $\varepsilon_s(A)$ to a diagram $D$ are identified with
morphisms from $A$ to $D(s)$.

\LayerTarget{tohoku:def:cogenerator-family}
We leave to the reader the task of developing the dual notion of a
\textit{family of cogenerators} of an abelian category. One can show that if
$\mathbf{C}$ is an abelian category satisfying axiom AB 5) (\textit{cf.} 1.5), then the
existence of a generator implies the existence of a cogenerator; we shall not
use this result. Thus the category of unital left modules over a ring $U$ with
identity always has a cogenerator. For example, if $U=\mathbf Z$, one may take
the group of rational numbers modulo $\mathbf Z$, or the torus
$\mathbf T=\mathbf R/\mathbf Z$.

\LayerTarget{tohoku:section:1.10}
\textbf{1.10. Injective and projective objects.}
\LayerTarget{tohoku:def:injective-object}
Recall that an object $M$ of an abelian category $\mathbf{C}$ is called
\textit{injective} if the functor $A\to\mathrm{Hom}(A,M)$, which is in
any case left exact, is exact; that is, if every morphism $u$ to $M$ from a
subobject $B$ of an object $A\in \mathbf{C}$ extends to a morphism from $A$ to $M$.
\LayerTarget{tohoku:def:injective-effacement}
A morphism $A\to M$ is called an \textit{injective effacement} of
$A$ if it is a monomorphism and if, for every monomorphism
$B\to C$ and every morphism $B\to A$, there is a
morphism $C\to M$ making the diagram commute:

\LayerTarget{tohoku:diagram:injective-effacement}
\[
\begin{array}{ccc}
B & \longrightarrow & C \\
\big\downarrow && \big\downarrow \\
A & \longrightarrow & M
\end{array}
\]
Thus the identity map of $M$ is an injective effacement if and only if $M$ is
injective; every monomorphism into an injective object is an injective
effacement.

\LayerTarget{tohoku:thm:1.10.1:p1}
\textsc{Theorem 1.10.1.} \textit{If $\mathbf{C}$ satisfies axiom \textup{AB 5)}
\textup{(}\textit{cf.} \textup{1.5)} and has a generator \textup{(}\textit{cf.} \textup{1.9)}, then for every $A\in \mathbf{C}$ there is a
monomorphism from $A$ into an injective object $M$.}

We shall in fact construct a functor
$M:A\to M(A)$, not additive in general,

\LayerTarget{tohoku:footnote:2}
\SourceFootnote{2}{The category $\mathbf{C}^\Sigma$ of inductive systems in $\mathbf{C}$ constructed on an
\textit{ordered} index set $I$ falls under this case. Indeed, in Example 1.7.h)
it suffices to consider the arrows $(i,j)$ with $i<j$ and the commutation
relations $(i,j)(j,k)=(i,k)$, in which the $e_s$ no longer occur.}

\TohokuContinuation{136}
